\section*{Bab \ref{ch:g10:storms-and-cyclones} --- Badai dan Siklon Tropis}

\begin{solution}{exo:g10:storms-and-cyclones:1}
\omterm{def:g10:storms-and-cyclones:thunderstorm}{Badai guntur} adalah badai yang dihasilkan oleh kumulonimbus, disertai \omterm{def:g10:storms-and-cyclones:lightning}{petir}
dan guntur, biasanya hujan deras dan embusan angin, kadang-kadang hujan es.
\omterm{def:g10:storms-and-cyclones:lightning}{Petir} adalah percikan listrik raksasa yang melepaskan muatan daerah-daerah
bermuatan dalam badai. Guntur adalah bunyi udara di saluran \omterm{def:g10:storms-and-cyclones:lightning}{petir}, yang
dipanaskan begitu tiba-tiba sehingga memuai dengan sentakan.
\end{solution}

\begin{solution}{exo:g10:storms-and-cyclones:2}
\omterm{def:g10:storms-and-cyclones:tropical-cyclone}{Siklon tropis} adalah badai yang lebarnya ratusan kilometer, terbentuk di atas
laut tropis yang hangat di sekitar pusat bertekanan rendah, dengan angin yang
berpilin masuk dan pita-pita \omterm{def:g10:storms-and-cyclones:thunderstorm}{badai guntur}. \omterm{def:g10:storms-and-cyclones:tropical-cyclone}{Mata} adalah pusatnya yang tenang
dan hampir tanpa awan; \omterm{def:g10:storms-and-cyclones:tropical-cyclone}{dinding mata} adalah cincin \omterm{def:g10:storms-and-cyclones:thunderstorm}{badai guntur} yang menjulang
di sekeliling \omterm{def:g10:storms-and-cyclones:tropical-cyclone}{mata}. Angin paling kuat di \omterm{def:g10:storms-and-cyclones:tropical-cyclone}{dinding mata}.
\end{solution}

\begin{solution}{exo:g10:storms-and-cyclones:3}
\num{70}~knot: kategori 1. \num{120}~knot: kategori 4.
\num{140}~knot: kategori 5. Dua yang terakhir adalah hurikan besar (kategori 3
atau lebih).
\end{solution}

\begin{solution}{exo:g10:storms-and-cyclones:4}
$343 \times 9 \approx \qty{3.1}{km}$, lalu $343 \times 15 \approx
\qty{5.1}{km}$. Kilatan kedua lebih jauh: badai itu sedang menjauh.
\end{solution}

\begin{solution}{exo:g10:storms-and-cyclones:5}
\omterm{def:g10:storms-and-cyclones:lightning}{Petir} cenderung menyambar titik tertinggi di sekitarnya: pohon yang berdiri
sendirian di lapangan adalah titik itu, dan arus sambaran mengalir menuruni
pohon dan menyebar melalui tanah di sekitarnya. Masuklah ke dalam bangunan
atau mobil yang tertutup, dan tetaplah di sana sampai lama setelah guntur
terakhir.
\end{solution}

\begin{solution}{exo:g10:storms-and-cyclones:6}
Satu knot adalah $1852/3600 \approx \qty{0.514}{m/s}$. \num{64}~knot: $64 \times
0.514 \approx \qty{32.9}{m/s}$, dan $32.9 \times 3.6 \approx
\qty{119}{km/h}$. \num{137}~knot: $\qty{70.5}{m/s}$, sekitar
\qty{254}{km/h} (tabel resmi, yang dibulatkan melalui satuan lain, memberikan
\qty{252}{km/h}).
\end{solution}

\begin{solution}{exo:g10:storms-and-cyclones:7}
$44 \times 86\,400 \approx \num{3.8e6}$ kilatan sehari, dan $3.8\times10^6
\times 365.25 \approx \num{1.4e9}$, sekitar 1.4 miliar setahun.
\end{solution}

\begin{solution}{exo:g10:storms-and-cyclones:8}
$\Delta p = 1010 - 935 = \qty{75}{hPa}$: gembungan sekitar \qty{75}{cm}.
\end{solution}

\begin{solution}{exo:g10:storms-and-cyclones:9}
\omterm{def:g10:storms-and-cyclones:tropical-cyclone}{Siklon tropis} terbentuk di atas samudra tropis yang hangat, sebagian besar di
antara lintang \qty{10}{\degree} dan \qty{30}{\degree} di kedua belahan bumi,
dengan sangat sedikit di Atlantik Selatan. Siklon itu hanyut ke barat, lalu
membelok menjauh ke arah kutub, tempat siklon itu melemah di atas air yang
lebih dingin. \omterm{def:g10:storms-and-cyclones:tropical-cyclone}{Siklon tropis} tidak ada di khatulistiwa karena di sana rotasi
Bumi tidak membelokkan udara yang mengalir masuk: tanpa belokan itu, udara
mengalir lurus ke dalam tekanan rendah dan mengisinya, sehingga tidak ada badai
berputar yang dapat dimulai.
\end{solution}

\begin{solution}{exo:g10:storms-and-cyclones:10}
$1031 - 1.56 \times 70 \approx \qty{922}{hPa}$; $1031 - 1.56 \times 40
\approx \qty{969}{hPa}$.
\end{solution}

\begin{solution}{exo:g10:storms-and-cyclones:11}
Di atas daratan, badai kehilangan pasokan air yang menguap, dan bersamanya
\omterm{def:g10:earth-energy-budget:latent-heat}{kalor laten} yang menggerakkannya; permukaan yang kasar juga memperlambat
angin. Di atas air yang dingin penguapan jauh lebih lemah, karena udara dingin
menampung uap air jauh lebih sedikit pada kejenuhan: pasokan \omterm{def:g10:earth-energy-budget:latent-heat}{kalor laten}
menyusut.
\end{solution}

\begin{solution}{exo:g10:storms-and-cyclones:12}
Pada \qty{0}{\degreeCelsius}: $\sqrt{1.4 \times 8.314 \times 273.15/0.02897}
\approx \qty{331}{m/s}$; pada \qty{30}{\degreeCelsius}, dengan $T =
\qty{303.15}{K}$: \qty{349}{m/s}. Untuk \qty{9}{s}: $9 \times 331 \approx
\qty{2980}{m}$ dan $9 \times 349 \approx \qty{3140}{m}$, selisihnya sekitar
\qty{160}{m}, sekitar \qty{5}{\percent}. Aturan itu memberikan \qty{3}{km}:
cukup baik untuk mengetahui seberapa dekat badai itu, dan memang hanya untuk
itulah aturan itu dibuat.
\end{solution}

\begin{solution}{exo:g10:storms-and-cyclones:13}
$\Delta p = 1010 - 940 = \qty{70}{hPa}$: gembungannya sekitar \qty{0.7}{m},
bagian angin $4.0 - 0.7 = \qty{3.3}{m}$, lebih dari empat kali lebih besar.
Saat pasang naik air mencapai $4.0 + 1.5 = \qty{5.5}{m}$ di atas muka laut
rata-rata.
\end{solution}

\begin{solution}{exo:g10:storms-and-cyclones:14}
\omterm{def:g10:storms-and-cyclones:tornado}{Tornado}: satu putaran adalah $2\pi \times 100 \approx \qty{628}{m}$, ditempuh
dalam $628/90 \approx \qty{7}{s}$. Hurikan: $2\pi \times 30\,000 \approx
\qty{188}{km}$, ditempuh dalam $188\,000/60 \approx \qty{3140}{s}$, sekitar
\qty{52}{min}. \omterm{def:g10:storms-and-cyclones:tornado}{Tornado} berputar sekitar 450 kali lebih cepat: \omterm{def:g10:storms-and-cyclones:tornado}{tornado} adalah
pusaran kecil yang sangat cepat, hurikan pusaran raksasa yang lebih lambat.
\end{solution}

\begin{solution}{exo:g10:storms-and-cyclones:15}
$24/3888 \approx \qty{0.6}{\percent}$; $414/3888 \approx
\qty{10.6}{\percent}$; $121/3888 \approx \qty{3.1}{\percent}$. Jika lintang
tidak berpengaruh, sekitar \qty{9}{\percent} badai, sekitar 350, akan bermula
dalam jarak \qty{5}{\degree} dari khatulistiwa; kenyataannya hanya 24, sekitar
lima belas kali lebih sedikit. Di dekat khatulistiwa rotasi Bumi terlalu lemah
untuk memulai putaran.
\end{solution}

\begin{solution}{pb:g10:storms-and-cyclones:1}
\textbf{1.} $\pi \times (665\,000)^2 \approx \qty{1.39e12}{m^2}$.

\textbf{2.} $0.015 \times 1.39\times10^{12} \approx \qty{2.08e10}{m^3}$.

\textbf{3.} Satu kilometer kubik adalah \qty{e9}{m^3}: sekitar \qty{21}{km^3}.

\textbf{4.} $2.08\times10^{10} \times 1000 \approx \qty{2.08e13}{kg}$.

\textbf{5.} $2.08\times10^{13} \times 2.5\times10^6 \approx \qty{5.2e19}{J}$.

\textbf{6.} $5.2\times10^{19}/86\,400 \approx \qty{6.0e14}{W}$.

\textbf{7.} $6.0\times10^{14}/1.22\times10^{17} \approx 0.005$, sekitar
\qty{0.5}{\percent}.

\textbf{8.} $240 \times 1.39\times10^{12} \approx \qty{3.3e14}{W}$. Badai itu
melepaskan kalor hampir dua kali cahaya Matahari yang diserap oleh seluruh
luasnya.

\textbf{9.} Badai itu tidak berjalan dengan cahaya Matahari yang jatuh di
atasnya: badai itu melepaskan kalor yang tersimpan sebagai uap air, yang
diuapkan dari laut hangat yang lebih luas dan ditarik masuk oleh angin. Badai
itu memusatkan energi yang dikumpulkan di atas daerah yang jauh lebih luas.

\textbf{10.} Massa yang sama dengan hujannya, sekitar \qty{2.1e13}{kg} sehari
(uap yang dibawa masuk harus diperbarui secepat uap itu tercurah sebagai
hujan).

\textbf{11.} $6.0\times10^{14}/82 \approx \qty{7.3e12}{m^2}$, sekitar lima kali
luas badai: badai itu harus menarik uap dari daerah yang jauh lebih luas, dan
anginnya membuat laut di bawahnya menguap jauh lebih cepat daripada rata-rata.

\textbf{12.} \qty{24.1}{g/kg} pada \qty{28}{\degreeCelsius} dibandingkan
dengan \qty{14.7}{g/kg} pada \qty{20}{\degreeCelsius}:
$24.1/14.7 \approx 1.64$, sekitar \qty{64}{\percent} lebih banyak.

\textbf{13.} Air hangat menguap dengan cepat dan udara hangat menampung uap
jauh lebih banyak: udara yang mengalir masuk membawa jauh lebih banyak \omterm{def:g10:earth-energy-budget:latent-heat}{kalor
laten} untuk dilepaskan, dan, dalam keadaan jenuh, udara itu naik dengan bebas.
Di atas laut yang lebih sejuk, pasokannya terlalu kecil untuk menggerakkan
mesin itu.

\textbf{14.} Di atas daratan tidak ada laut yang dapat menguap: pasokan \omterm{def:g10:earth-energy-budget:latent-heat}{kalor
laten} terhenti, dan tanah yang kasar memperlambat angin.

\textbf{15.} $1.5\times10^{12} \times 86\,400 \approx \qty{1.3e17}{J}$.

\textbf{16.} $96 \times 1852/3600 \approx \qty{49.4}{m/s}$.

\textbf{17.} $\frac12 \times 1.16 \times 49.4^2 \approx \qty{1.4e3}{J}$.

\textbf{18.} $5.2\times10^{19}/1.3\times10^{17} = 400$: \omterm{def:g10:earth-energy-budget:latent-heat}{kalor laten} itu
sekitar 400 kali energi yang diberikan kepada angin.

\textbf{19.} Sebagian besar kalor itu menghangatkan udara atas badai, yang
menyebar di puncaknya dan melepaskan kalor ke angkasa sebagai radiasi
inframerah; angin itu sendiri kehilangan energinya karena gesekan dengan laut
dan tanah.

\textbf{20.} $1/400 = \qty{0.25}{\percent}$ (skrip memberikan rasio
402): \textbf{sekitar satu bagian dari 400 kalor yang dilepaskan oleh
pengembunan menjadi angin}.
\end{solution}
